\section{Entity Matching}
\label{sec:entity_matching}
Entity matching identifies records across different datasets that describe the same real-world entity~\cite{christen2012data, peeters2024entity}. This section compares the data engineers' manually configured entity matchers against ML-based matchers trained on human-labeled and LLM-labeled data.

\textbf{Human Pipeline.}\label{sec:em_human} The data engineers manually configured blocking methods, attribute comparators, and matching thresholds for each dataset pair using PyDI. Blocking keys are derived from entity names (e.g., longest or first token) to prune the comparison space. PyDI provides typed comparators (string, date, numeric) that the data engineers selected for each attribute. Depending on the dataset pair, the data engineers used either rule-based matchers employing attribute weights and a similarity threshold or ML classifiers trained on human-labeled pairs.

\textbf{LLM Pipeline.}
\label{sec:em_approach}
The LLM pipeline replaces both human labeling and manual matcher configuration with LLM-generated training data and automated model selection. We employ an active learning workflow to select entity pairs to be labeled by the LLM. We use embedding-based blocking to generate a pool of candidate pairs: records are embedded using OpenAI's text-embedding-3-small model and the $k{=}20$ nearest neighbors from each other dataset are retrieved. GPT-5.2 labels pairs from this pool in order of decreasing similarity, stopping for each query entity once at least one positive and two negative matches have been found. Two candidates from the bottom of the similarity ranking are additionally labeled per query to increase diversity. The target number of seed pairs is configurable; the process stops once this target is reached or a labeling budget is exhausted, yielding approximately 100 labeled seed pairs per dataset pair in our setup (see footnote~\ref{fn:prompts} for all prompts).

The seed pairs are used to initialize an active learning loop. Multiple classifiers (RandomForest, XGBoost, GradientBoosting, HistGradientBoosting, LogisticRegression) are trained on string similarity features (Jaccard, Jaro-Winkler, Levenshtein, cosine, and others) computed for each candidate pair, using the same datatype-dependent PyDI similarity metrics as the human pipeline. Hyperparameters are tuned via randomized search. The active learning loop identifies the pairs in the candidate pool on which the classifiers disagree most, measured by the variance of their prediction scores, and selects 100 such candidates for GPT-5.2 to label per iteration. Newly labeled pairs are added to the training set and all classifiers are retrained. This repeats until the training set reaches a target size or the labeling budget is exhausted. The final training set is augmented with 20\% randomly sampled pairs from the pool to counteract the selection bias toward difficult comparisons. The best combination of training set variant, classifier, and threshold is selected per dataset pair by F1 on a GPT-5.2-labeled validation set. Across all dataset pairs, RandomForest and XGBoost are the classifiers most frequently selected.


% Entity matching results: Human Config vs Human Labels vs LLM Labels

\begin{table}[t]
\caption{Entity matching F1 scores on held-out test sets. \emph{Human Config}: the data engineers' pipeline. \emph{Human Labels} and \emph{LLM Labels}: identical ML matchers trained on human-labeled vs.\ LLM-labeled training data.}
\label{tab:em_results}
\centering
\footnotesize
\begin{tabular}{@{}ll ccc@{}}
\toprule
 & & \textbf{Human} & \textbf{Human} & \textbf{LLM} \\
\textbf{Use Case} & \textbf{Dataset Pair} & \textbf{Config} & \textbf{Labels} & \textbf{Labels} \\
\midrule
\multirow{2}{*}{Games}
 & DBpedia--Metacritic & .930 & .826 & .849 \\
 & DBpedia--Sales & .927 & .839 & .979 \\
\midrule
\multirow{2}{*}{Companies}
 & DBpedia--Forbes & .857 & .954 & .939 \\
 & Forbes--FullContact & .870 & .898 & .897 \\
\midrule
\multirow{2}{*}{Music}
 & Discogs--MusicBrainz & .800 & .991 & .990 \\
 & Last.fm--MusicBrainz & .977 & .988 & .968 \\
\midrule
\multicolumn{2}{@{}l}{\textit{Average}} & .894 & .916 & .937 \\
\bottomrule
\end{tabular}
\end{table}


\textbf{Results.}
\label{sec:em_results}
Table~\ref{tab:em_results} compares three configurations on held-out test sets ranging from 139 to 1,000 manually labeled pairs per dataset pair: The data engineers' manually configured pipeline (Human Config), ML matchers trained on human-labeled data (Human Labels), and the same ML matchers trained on LLM-labeled data (LLM Labels). The Human Labels and LLM Labels columns isolate the effect of label quality, as both use identical feature generation and model selection.

The LLM-labeled matchers achieve the highest average F1 of .937, followed by human-labeled matchers at .916 and the manually configured pipeline at .894. The manually configured matchers are competitive in Games (F1 .930 and .927) and reach the highest individual score in Music Last.fm--MusicBrainz (.977), but perform substantially worse on Music Discogs--MusicBrainz (.800) and Companies (.857 and .870). 

Comparing Human Labels and LLM Labels, performance is nearly identical in Music (F1 $>$ .96 for both) and Companies (within 1.5 percentage points). The largest gain from LLM labeling appears in Games DBpedia--Sales (+14.0\%), possibly due to labeling errors in the manually created training set or because active learning selected more relevant examples compared to the data engineers' selection. The active learning variant achieved the best F1 in five of six cases.



\textbf{Cluster Post-Processing.} False positive correspondences can link unrelated entities into oversized clusters containing more records than the number of source datasets. In order to split such clusters, we remove potentially wrong correspondences by applying maximum bipartite filtering. After cleaning, clusters contain at most three records (one per source), dominated by pairs (62--79\%).


\textbf{Discussion.}
\label{sec:em_limitations}
The results show that ML matchers trained on LLM-generated labels match or exceed the effectiveness of both manually configured matchers and matchers trained on human-labeled data. The effectiveness of the LLM-generated training sets likely results from both the active learning workflow selecting informative pairs and GPT-5.2's ability to correctly label them. Embedding-based blocking with $k{=}20$ neighbors generates a large candidate space (up to 837k pairs across all dataset pairs). Labeling all candidates directly with the LLM would cost approximately \$466, whereas the active learning approach labels only 600--2,658 examples per dataset pair at an overall cost of \$4.95 for all three case studies (Table~\ref{tab:pipeline_costs}), making the active learning-based approach two orders of magnitude cheaper. 